\pagebreak

\section{Clock Manager and Clock Buffer Library File - xilinx/clocks.3pl}

    \index{library!xilinx/clocks.3pl}
    This library creates any clock managers or clock buffers for all Xilinx
    FPGAs.
    
    This library must be incorporated by using a {\bf module} statement -
    \begin{lstlisting}[gobble=4, language=threepl]
       module "xilinx/clocks.3pl"
    \end{lstlisting}
    
    {\bf This library is already incorporated by all FPGA family header files.}
    
    The procedures in this library attempt to propagate clock frequencies and
    periods from inputs to outputs using the parameters controlling clock
    division or synthesis. For that reason most take an error exit if an input
    clock does not have a frequency-period attribute.

    At the moment clkmmcm(), pll\_adv() are exceptions to the last paragraph.
    They will be changed to conform ASAP.
    
    Where clock input pads, clock buffers and clock manager are interconnected
    clock mode variables are used (required by parameter modes). Where a clock
    is derived from CLB logic the inbuilt function {\bf clockfromexpr()}
    \autorefph{sec:ipclockfromexpr} can be used to assign a value mode variable
    to a clock mode variable. Where input delays are required for an input clock
    path (see {\bf idelay()} \autorefph{sec:lpidelay}) the input should be input
    mode rather than clock mode. The output from the delay can then be convert
    to clock mode by using {\bf clockfromexpr()}.



    \subsection{library modules}
    
        None.


    \subsection{library procedures}

        \subsubsection{clkdcm - create a Spartan3, Virtex2, Virtex2Pro, Virtex4 or Virtex5 DCM digital clock manager}\label{sec:lpclkdcm}
            
            {\bf Library: xilinx/clocks.3pl}

            {\bf Prototype:} \\
            \vsp
            \begin{lstlisting}[gobble=12, language=threepl]
               global procedure
               clkdcm (
                   clock({clk0, clk90, clk180, clk270, clk2x, clk2x180, clkdv, clkfx, clkfx180}),
                   value({locked, psdone, drdy}, "log"), value(dout, "uint:16"), value(status, "uint:8")
                   <-
                   str(loc), log({clkin_divide_by_2, duty_cycle_correction, startup_wait}),
                   str(clkdv_divide),
                   float(clkin_period), int({clkfx_divide, clkfx_multiply, phase_shift}),
                   str({clk_feedback,
                   clkout_phase_shift, dfs_frequency_mode, dll_frequency_mode, deskew_adjust}),
                   clock({clkin, clkfb, psclk}), value({reset, psincdec, psen, dwe, den}, "log"),
                   value(din, "bits:16"), value(daddr, "uint:7")
               ) {}
            \end{lstlisting}

            {\bf Output parameters:} \\
            \begin{tabular}{| l | l | l | l |}
            \hline
                           & mode  & type   & \\
            \hline \hline
            {\bf clk0}     & clock &         & output clock \\ \hline
            {\bf clk90}    & clock &         & output clock shifted 90$\deg$ \\ \hline
            {\bf clk180}   & clock &         & output clock shifted 180$\deg$ \\ \hline
            {\bf clk270}   & clock &         & output clock shifted 270$\deg$ \\ \hline
            {\bf clk2x}    & clock &         & output clock at double frequency \\ \hline
            {\bf clk2x180} & clock &         & output clock at double frequency shifted 180$\deg$ \\ \hline
            {\bf clkdv}    & clock &         & divided output clock \\ \hline
            {\bf clkfx}    & clock &         & frequency synthesised output clock \\ \hline
            {\bf clkfx180} & clock &         & frequency synthesised output clock shifted 180$\deg$ \\ \hline
            {\bf locked}   & value & log     & DLL is locked \\ \hline
            {\bf status}   & value & uint:8  & status \\ \hline
            {\bf psdone}   & value & log     & phase shift has finished \\ \hline
            {\bf drdy}     & value & log     & DCM ready for reconfiguration \\ \hline
            {\bf dout}     & value & bits:16 & status and reconfiguration data output \\ \hline
            \end{tabular}

            {\bf Input parameters:} \\            
            \begin{tabular}{| l | l | l | l |}
            \hline
                           & mode  & type   & \\
            \hline \hline
            {\bf loc}                     & immediate & str   & optional location \\ \hline
            {\bf clkin\_divide\_by\_2}    & immed     & log   & DLL clock output at $f_{in}/2$ \\ \hline
            {\bf duty\_cycle\_correction} & immediate & log   & correct duty cycle to 50\%  \\ \hline
            {\bf startup\_wait}           & immediate & log   & delay startup on lock \\ \hline
            {\bf clkdv\_divide}           & immed     & str   & $n$ for clock output at $f_{in}/n$ \\ \hline
            {\bf clkin\_period}           & immediate & float & input clock period, ns \\ \hline
            {\bf clkfx\_divide}           & immed     & int   & $m$ for DFS clock output at $f_{in} \times n / m$ \\ \hline
            {\bf clkfx\_multiply}         & immed     & int   & $n$ for DFS clock output at $f_{in} \times n / m$ \\ \hline
            {\bf phase\_shift}            & immediate & int   & fixed phase shift value \\ \hline
            {\bf clk\_feedback}           & immediate & str   & "1x", "2x", "none" \\ \hline
            {\bf clkout\_phase\_shift}    & immediate & str   & "none", "fixed", "variable" \\ \hline
            {\bf dfs\_frequency\_mode}    & immediate & str   & "high" or "low" \\ \hline
            {\bf dll\_frequency\_mode}    & immediate & str   & "high" or "low" \\ \hline
            {\bf deskew\_adjust}          & immediate & str   & "system\_synchronous", "source\_synchronous", or 0 - 15 \\ \hline
            {\bf clkin}                   & clock     &       & input clock \\ \hline
            {\bf clkfb}                   & clock     &       & clock feedback \\ \hline
            {\bf psclk}                   & clock     &       & phase shift clock\\ \hline
            {\bf reset}                   & value     & log   & reset \\ \hline
            {\bf psincdec}                & value     & log   & phase shift increment or decrement\\ \hline
            {\bf psen}                    & value     & log   & phase shift enable\\ \hline
            {\bf dwe}                     & value     & log   & reconfiguration data write enable\\ \hline
            {\bf den}                     & value     & log   & select reconfiguration data or status output\\ \hline
            {\bf din}                     & value     & log   & reconfiguration data input \\
            {\bf daddr}                   & value     & log   & reconfiguration address input \\ \hline
            \end{tabular}

            {\bf In-line target execution:} \\
            no

            {\bf description:} \\
            Creates a Virtex4 or Virtex5 DCM digital clock manager. See Xilinx
            documentation for signal and attribute details. The first group of input
            parameters are immediate mode and are attributes. The remaining input
            parameters are input signals, clock mode for clocks and value mode for
            control signals.

            \heading{Attributes}

                {\bf loc} is an optional location string from the
                following list - "DLL0P", "DLL1P", "DLL2P", "DLL3P", "DLL0S", "DLL1S",
                "DLL2S", "DLL3S". These select from the four primary and four secondary
                delay-locked loops. If "" it is ignored.

                {\bf clkin\_divide\_by\_2} - see Xilinx documentation.

                If {\bf duty\_cycle\_correction} is true or missing then duty cycle
                correction is enabled, giving a DLL output duty cycle of 50\%. If
                this attribute is false then the duty cycle will be the same as that
                of the input.

                {\bf startup\_wait} if true causes startup following configuration to
                be delayed until the DCM achieves lock.

                {\bf clkfx\_divide} is the frequency divider for the DFS (1 to 32).

                {\bf clkfx\_multiply} is the frequency multiplier for the DFS (2 to 32).


                {\bf phase\_shift} is the phase shift value which must be in the
                range -255 to +255 inclusive. A unit value represents 1/256 of a cycle.

                {\bf clkdv\_divide} indicates that the divided
                frequency output is being used. The division ratio is a
                string argument which must be one of "1.5", "2.0", "2.5", "3.0",
                "3.5", "4.0", "4.5", "5.0", "5.5", "6.0", "6.5", "7.0", "7.5", "8.0",
                "9.0", "10.0", "11.0", "12.0", "13.0", "14.0", "15.0" or "16.0".

                {\bf clkin\_period} is the input clock period. This must be
                specified if the DFS is used.

                {\bf clk\_feedback} must be "2X" if the clock feedback is at twice
                input frequency. By default the feedback is at the input frequency, or
                this attribute can be explcitly "1X" to achieve the same.

                {\bf clkout\_phase\_shift} is a string selecting the mode of the
                phase shifter. If omitted or "none" the outputs are not phase shifted.
                If this argument is "fixed" a clock skew is added to the input clock
                shifting the phase of all clock outputs. The phase shift is specified by
                the next input argument. If "variable" is specified the phase shift
                is initially set as specified by the next argument, but it can be
                dynamically incremented or decremented - see the 17th to 19th input
                arguments. On reset the phase shift returns to the initial value.

                {\bf dfs\_frequency\_mode} indicates that the DFS is on
                the high frequency range. The low and high frequency ranges are
                given in the data sheets.

                {\bf dll\_frequency\_mode} indicates that the DLL is in
                the high frequency range (see data sheets). In this mode the 90 degree
                and 270 degree outputs and the 2X frequency outputs  are not available.

                {\bf deskew\_adjust} controls the amount of delay in the feedback
                path. Its values can be "system\_synchronous" (default),
                "source\_synchronous" or "0 to 15".

                {\bf dcm\_performance\_mode} optimises the DCM for high frequency low jitter
                ("max\_speed", the default) or low frequency and maximum phase shift range
                ("max\_range").

            \heading{input signals}
            
                {\bf din}, {\bf daddr}, {\bf dwe},  {\bf den} and {\bf dclk}
                are only available on Virtex4 and Virtex5.

                {\bf clkin} is the input clock. This argument must be
                supplied and is the first target mode input argument.

                {\bf clkfb} is the clock feedback. This argument must be
                supplied. This is usually a buffered version of the output clock at
                phase 0.

                {\bf reset} is an optional reset. This resets the
                DLL, DFS and phase shifter and re-establishes lock.

                {\bf psincdec} is true if a phase increment is to be
                executed or false if a phase decrement is to be executed. Execution is
                controlled by the following two inputs.

                {\bf psen} is the phase increment/decrement enable.
                This must be held true for one period of the phase shift adjustment
                clock (the next input argument). 

                {\bf psclk} is the phase increment/decrement clock.

                {\bf din} is data input for a dynamic reconfiguration write.

                {\bf daddr} is the address input for a dynamic reconfiguration
                write.

                {\bf dwe} is driven high for one {\bf dclk} cycle along with
                {\bf den} to execute a dynamic reconfiguration write.

                {\bf den} is driven high for one {\bf dclk} cycle to execute
                a dynamic reconfiguration read or write.

                {\bf dclk} is a clock controlling dynamic reconfiguration
                operations.

            \heading{output signals}

                {\bf drdy} and {\bf dout} are only available on Virtex4 and
                Virtex5.
                
                {\bf status} is only available on Spartan3, Virtex2 and Virtex2Pro.
                
                {\bf clk0}, {\bf clk90}, {\bf clk180} and {\bf clk270} are clock
                outputs at the input frequency. These are at phases 0, 90, 180 and 270
                degree respectively relative to the input clock.

                {\bf clk2x} and {\bf clk2x180} are at twice the input frequency
                at phases 0 and 180 degree respectively.

                {\bf clkdv} is the divided clock frequency where the division ratio
                is given by the 3rd input argument.

                {\bf clkfx} and {\bf clkfx180} are the frequency synthesised clock
                at phases 0 and 180 degree respectively. The frequency is the input
                multiplied by the value of {\bf clkfx\_multiply} and divided by the value
                of {\bf clkfx\_divide}.

                {\bf locked} is high when the DLL and DFS are locked to the input
                clock. This signal has a null clock domain.

                {\bf status} is an 8 bit status (see Xilinx unified library
                documentation). This signal has a null clock domain.

                {\bf psdone} goes high when a phase shift adjustment has been done.
                It remains high for one period of the phase adjustment clock. Note
                that this signal has a null clock domain, but it is synchronous
                with {\bf psclk}.

                {\bf drdy} goes high when the dynamic configuration port finishes
                an operation. On a read the data on {\bf do} must be captured on
                the next rising {\bf dclk} edge. When writing the next operation
                can commence on the same clock cycle in which {\bf drdy} goes high.

                {\bf dout} is the data output when reading dynamic configuration data.
                The data on {\bf do} must be captured on
                the rising {\bf dclk} edge when {\bf drdy} is high.



        \subsubsection{clkdll - create a Spartan2, Virtex or VirtexE DLL delay-locked loop clock manager}\label{sec:lpclkdll}
            
            {\bf Library: xilinx/clocks.3pl}

            {\bf Prototype:} \\
            \vsp
            \begin{lstlisting}[gobble=12, language=threepl]
               global procedure
               clkdlle (
                   clock({clk0, clk90, clk180, clk270, clk2x, clkdv}),
                   value(locked, "log")
                   <-
                   str(loc), log({duty_cycle_correction, startup_wait}),
                   int(clkdv_divide),
                   clock({clkin, clkfb}), value(reset, "log")
               ) {}
            \end{lstlisting}

            {\bf Output parameters:} \\
            \begin{tabular}{| l | l | l | l |}
            \hline
                           & mode  & type   & \\
            \hline \hline
            {\bf clk0}     & clock &        & output clock \\ \hline
            {\bf clk90}    & clock &        & output clock shifted 90$\deg$ \\ \hline
            {\bf clk180}   & clock &        & output clock shifted 180$\deg$ \\ \hline
            {\bf clk270}   & clock &        & output clock shifted 270$\deg$ \\ \hline
            {\bf clk2x}    & clock &        & output clock at double frequency \\ \hline
            {\bf clkdv}    & clock &        & divided output clock \\ \hline
            \end{tabular}

            {\bf Input parameters:} \\            
            \begin{tabular}{| l | l | l | l |}
            \hline
                           & mode  & type   & \\
            \hline \hline
            {\bf loc}                     & immediate & str   & optional location \\ \hline
            {\bf clkin\_divide\_by\_2}    & immediate & log   & divide clock by 2 \\ \hline
            {\bf duty\_cycle\_correction} & immediate & log   & correct duty cycle \\ \hline
            {\bf startup\_wait}           & immediate & log   & delay startup on lock \\ \hline
            {\bf clkdv\_divide}           & immediate & int   & clock division ratio \\ \hline
            {\bf clkin}                   & clock     &       & input clock \\ \hline
            {\bf clkfb}                   & clock     &       & clock feedback \\ \hline
            {\bf reset}                   & value     & log   & reset \\ \hline
            \end{tabular}

            {\bf In-line target execution:} \\
            no

            {\bf description:} \\
            Creates a Virtex DLL delay-locked loop clock manager. See Xilinx
            documentation for signal and attribute details. The first group of input
            parameters are immediate mode and are attributes. The remaining input
            parameters are input signals, clock mode for clocks and value mode for
            control signals.

            \heading{attributes}

                {\bf loc} is an optional location string from the
                following list - "DLL0P", "DLL1P", "DLL2P", "DLL3P", "DLL0S", "DLL1S",
                "DLL2S", "DLL3S". These select from the four primary and four secondary
                delay-locked loops. If "" it is ignored.

                {\bf clkin\_divide\_by\_2} - see Xilinx documentation.

                If {\bf duty\_cycle\_correction} is true or missing then duty cycle
                correction is enabled, giving a DLL output duty cycle of 50\%. If
                this attribute is false then the duty cycle will be the same as that
                of the input.

                {\bf startup\_wait} if true causes startup following configuration to
                be delayed until the DLL achieves lock.

                {\bf clkdv\_divide} indicates that the divided
                frequency output is being used. The division ratio is a
                string argument which must be an integer from 1 to 32 inclusive.

            \heading{input signals}

                {\bf clkin} is the input clock. This argument must be
                supplied and is the first target mode input argument.

                {\bf clkfb} is the clock feedback. This argument must be
                supplied. This is usually a buffered version of the output clock at
                phase 0.

                {\bf reset} is an optional reset. This resets the
                DLL, DFS and phase shifter and re-establishes lock.
                operations.

            \heading{output signals}

                {\bf clk0}, {\bf clk90}, {\bf clk180} and {\bf clk270} are clock
                outputs at the input frequency. These are at phases 0, 90, 180 and 270
                degree respectively relative to the input clock.

                {\bf clk2x} is at twice the input frequency.

                {\bf clkdv} is the divided clock frequency where the division ratio
                is given by the 3rd input argument.

                {\bf locked} is high when the DLL and DFS are locked to the input
                clock. This signal has a null clock domain.






        \subsubsection{clkdlle - create an extended Spartan2, Virtex or VirtexE DLL delay-locked loop clock manager}\label{sec:lpclkdlle}
            
            {\bf Library: xilinx/clocks.3pl}

            {\bf Prototype:} \\
            \vsp
            \begin{lstlisting}[gobble=12, language=threepl]
               global procedure
               clkdlle (
                   clock({clk0, clk90, clk180, clk270, clk2x, clk2x180, clkdv}),
                   value(locked, "log")
                   <-
                   str(loc), log({duty_cycle_correction, startup_wait}),
                   int(clkdv_divide),
                   clock({clkin, clkfb}), value(reset, "log")
               ) {}
            \end{lstlisting}

            {\bf Output parameters:} \\
            \begin{tabular}{| l | l | l | l |}
            \hline
                           & mode  & type   & \\
            \hline \hline
            {\bf clk0}     & clock &        & output clock \\ \hline
            {\bf clk90}    & clock &        & output clock shifted 90$\deg$ \\ \hline
            {\bf clk180}   & clock &        & output clock shifted 180$\deg$ \\ \hline
            {\bf clk270}   & clock &        & output clock shifted 270$\deg$ \\ \hline
            {\bf clk2x}    & clock &        & output clock at double frequency \\ \hline
            {\bf clk2x180} & clock &        & output clock at double frequency shifted 180$\deg$ \\ \hline
            {\bf clkdv}    & clock &        & divided output clock \\ \hline
            \end{tabular}

            {\bf Input parameters:} \\            
            \begin{tabular}{| l | l | l | l |}
            \hline
                           & mode  & type   & \\
            \hline \hline
            {\bf loc}                     & immediate & str   & optional location \\ \hline
            {\bf clkin\_divide\_by\_2}    & immediate & log   & divide clock by 2 \\ \hline
            {\bf duty\_cycle\_correction} & immediate & log   & correct duty cycle \\ \hline
            {\bf startup\_wait}           & immediate & log   & delay startup on lock \\ \hline
            {\bf clkdv\_divide}           & immediate & int   & clock division ratio \\ \hline
            {\bf clkin}                   & clock     &       & input clock \\ \hline
            {\bf clkfb}                   & clock     &       & clock feedback \\ \hline
            {\bf reset}                   & value     & log   & reset \\ \hline
            \end{tabular}

            {\bf In-line target execution:} \\
            no

            {\bf description:} \\
            Creates a Virtex DLLE extended delay-locked loop clock manager. See
            Xilinx documentation for signal and attribute details. The first group
            of input parameters are immediate mode and are attributes. The remaining
            input parameters are input signals, clock mode for clocks and value mode
            for control signals.

            \heading{attributes}

                {\bf loc} is an optional location string from the
                following list - "DLL0P", "DLL1P", "DLL2P", "DLL3P", "DLL0S", "DLL1S",
                "DLL2S", "DLL3S". These select from the four primary and four secondary
                delay-locked loops. If "" it is ignored.

                {\bf clkin\_divide\_by\_2} - see Xilinx documentation.

                If {\bf duty\_cycle\_correction} is true or missing then duty cycle
                correction is enabled, giving a DLL output duty cycle of 50\%. If
                this attribute is false then the duty cycle will be the same as that
                of the input.

                {\bf startup\_wait} if true causes startup following configuration to
                be delayed until the DLL achieves lock.

                {\bf clkdv\_divide} indicates that the divided
                frequency output is being used. The division ratio is a
                string argument which must be an integer from 1 to 32 inclusive.

            \heading{input signals}

                {\bf clkin} is the input clock. This argument must be
                supplied and is the first target mode input argument.

                {\bf clkfb} is the clock feedback. This argument must be
                supplied. This is usually a buffered version of the output clock at
                phase 0.

                {\bf reset} is an optional reset. This resets the
                DLL, DFS and phase shifter and re-establishes lock.
                operations.

            \heading{output signals}

                {\bf clk0}, {\bf clk90}, {\bf clk180} and {\bf clk270} are clock
                outputs at the input frequency. These are at phases 0, 90, 180 and 270
                degree respectively relative to the input clock.

                {\bf clk2x} and {\bf clk2x180} are at twice the input frequency
                at phases 0 and 180 degree respectively.

                {\bf clkdv} is the divided clock frequency where the division ratio
                is given by the 3rd input argument.

                {\bf locked} is high when the DLL and DFS are locked to the input
                clock. This signal has a null clock domain.






        \subsubsection{clkdllhf - create a high frequency Spartan2, Virtex or VirtexE DLL delay-locked loop clock manager}\label{sec:lpclkdllhf}
            
            {\bf Library: xilinx/clocks.3pl}

            {\bf Prototype:} \\
            \vsp
            \begin{lstlisting}[gobble=12, language=threepl]
               global procedure
               clkdlle (
                   clock({clk0, clk180, clkdv}),
                   value(locked, "log")
                   <-
                   str(loc), log({duty_cycle_correction, startup_wait}),
                   int(clkdv_divide),
                   clock({clkin, clkfb}), value(reset, "log")
               ) {}
            \end{lstlisting}

            {\bf Output parameters:} \\
            \begin{tabular}{| l | l | l | l |}
            \hline
                           & mode  & type   & \\
            \hline \hline
            {\bf clk0}     & clock &        & output clock \\ \hline
            {\bf clk180}   & clock &        & output clock shifted 180$\deg$ \\ \hline
            {\bf clkdv}    & clock &        & divided output clock \\ \hline
            \end{tabular}

            {\bf Input parameters:} \\            
            \begin{tabular}{| l | l | l | l |}
            \hline
                           & mode  & type   & \\
            \hline \hline
            {\bf loc}                     & immediate & str   & optional location \\ \hline
            {\bf clkin\_divide\_by\_2}    & immediate & log   & divide clock by 2 \\ \hline
            {\bf duty\_cycle\_correction} & immediate & log   & correct duty cycle \\ \hline
            {\bf startup\_wait}           & immediate & log   & delay startup on lock \\ \hline
            {\bf clkdv\_divide}           & immediate & int   & clock division ratio \\ \hline
            {\bf clkin}                   & clock     &       & input clock \\ \hline
            {\bf clkfb}                   & clock     &       & clock feedback \\ \hline
            {\bf reset}                   & value     & log   & reset \\ \hline
            \end{tabular}

            {\bf In-line target execution:} \\
            no

            {\bf description:} \\
            Creates a Virtex DLL high frequency delay-locked loop clock manager. See
            Xilinx documentation for signal and attribute details. The first group
            of input parameters are immediate mode and are attributes. The remaining
            input parameters are input signals, clock mode for clocks and value mode
            for control signals.

            \heading{attributes}

                {\bf loc} is an optional location string from the
                following list - "DLL0P", "DLL1P", "DLL2P", "DLL3P", "DLL0S", "DLL1S",
                "DLL2S", "DLL3S". These select from the four primary and four secondary
                delay-locked loops. If "" it is ignored.

                {\bf clkin\_divide\_by\_2} - see Xilinx documentation.

                If {\bf duty\_cycle\_correction} is true or missing then duty cycle
                correction is enabled, giving a DLL output duty cycle of 50\%. If
                this attribute is false then the duty cycle will be the same as that
                of the input.

                {\bf startup\_wait} if true causes startup following configuration to
                be delayed until the DLL achieves lock.

                {\bf clkdv\_divide} indicates that the divided
                frequency output is being used. The division ratio is a
                string argument which must be an integer from 1 to 32 inclusive.

            \heading{input signals}

                {\bf clkin} is the input clock. This argument must be
                supplied and is the first target mode input argument.

                {\bf clkfb} is the clock feedback. This argument must be
                supplied. This is usually a buffered version of the output clock at
                phase 0.

                {\bf reset} is an optional reset. This resets the
                DLL, DFS and phase shifter and re-establishes lock.
                operations.

            \heading{output signals}

                {\bf clk0} and {\bf clk180} are clock
                outputs at the input frequency. These are at phases 0 and 180
                degree respectively relative to the input clock.

                {\bf clkdv} is the divided clock frequency where the division ratio
                is given by the 3rd input argument.

                {\bf locked} is high when the DLL and DFS are locked to the input
                clock. This signal has a null clock domain.






        \subsubsection{simpleclock - create a clock of specified frequency}\label{sec:lpsimpleclock}
            
            {\bf Library: xilinx/clocks.3pl}

            {\bf Prototype:} \\
            \vsp
            \begin{lstlisting}[gobble=12, language=threepl]
               global procedure
               simpleclock (ident(v), clock(in), float(freq), float(tol)) {}
            \end{lstlisting}

            {\bf Output parameters:} \\
            none.

            {\bf Input parameters:} \\            
            \begin{tabular}{| l | l | l | l |}
            \hline
                                          & mode  & type   & \\
            \hline \hline
            {\bf v}     & -         & -       & new clock variable identifier \\ \hline
            {\bf in}    & clock     & -       & input clock                   \\ \hline
            {\bf freq}  & immediate & float   & requested frequency           \\ \hline
            {\bf tol}   & immediate & float   & frequency tolerance           \\ \hline
            \end{tabular}

            {\bf In-line target execution:} \\
            no

            {\bf description:} \\
            Creates a clock variable using a clkdcm and clock buffer as its source.
            {\bf At present it can only be used on FPGA platforms that contain DCMs.}
            
            {\bf v} is an identifier for the new clock mode variable.
            
            {\bf in} is a clock that can be used to drive the DCM.
            
            {\bf freq} is the requested frequency in Hz. If this has no argument
            then the clock frequency will be the same as that of the input clock.
            This can be used to simply add a DCM and clock buffer to an input clock
            from an input buffer.

            {\bf tol} is the allowable tolerance on the requested frequency in Hz.

            If a frequency within the specified tolerance cannot be generated then
            a fatal error occurs.

            This procedure does not attempt to optimise the use of DCMs to generate
            multiple clocks. It uses a new DCM for each call.
            
            The procedure determines the multiplier and divisor which will produce
            the closest synthesiser frequency to that requested. If the requested
            frequency is lower than the input frequency the procedure also attempts
            to find a clock division ratio close to that required to achieve the
            requested frequency. The closest match from the above is then chosen. If
            the resulting frequency is not within tolerance a fatal error occurs.
            
            The clock variable will have the generated frequency attribute, not
            that of the requested frequency.
            
            Example of simple clock input.

            \begin{lstlisting}[gobble=12, language=threepl]
               iopins(c32_pin, "str", "P94", frequency=32.0e6, ibufg=true, iostandard=lvcmos33);

               clock(c32_in, loc=c32_pin);     // clock input from pin and input buffer
               simpleclock(global.c32, c32_in);// global clock
            \end{lstlisting}


            Example of multiple global clocks generated from c32 clock.

            \begin{lstlisting}[gobble=12, language=threepl]
               simpleclock(global.c16, c32, 16.0e6, 0.1);    // 16 MHz clock
               simpleclock(global.c25, c32, 25.0e6, 0.1);    // 25 MHz clock
            \end{lstlisting}




        \subsubsection{clkmmcm\_adv - create an advanced MMCM multi-mode clock manager }\label{sec:lpmmcma}
            
            {\bf Library: xilinx/clocks.3pl}

            {\bf Prototype:} \\
            \vsp
            \begin{lstlisting}[gobble=12, language=threepl]
               global procedure
               clkmmcm_adv (
                   clock({ clkout0, clkout1, clkout2, clkout3, clkout4, clkout5, clkout6,
                           clkout0b, clkout1b, clkout2b, clkout3b, clkfbout, clkfboutb}),
                   value({clkfbstopped, clkinstopped, drdy, locked, psdone}, "log"),
                   value(dout, "bits:16")
                   <-
                   str(loc),

                   log({   clkfbout_use_fine_ps, clkout0_use_fine_ps, clkout1_use_fine_ps,
                           clkout2_use_fine_ps, clkout3_use_fine_ps, clkout4_use_fine_ps,
                           clkout5_use_fine_ps, clkout6_use_fine_ps, clock_hold, startup_wait}),
                   int(divclk_divide),
                   float({ clkfbout_phase, clkfbout_mult_f,
                           clkout0_divide_f,
                           clkout0_duty_cycle, clkout0_phase, clkout1_divide_f,
                           clkout1_duty_cycle, clkout1_phase, clkout2_divide_f,
                           clkout2_duty_cycle, clkout2_phase, clkout3_divide_f,
                           clkout3_duty_cycle, clkout3_phase, clkout4_divide_f,
                           clkout4_duty_cycle, clkout4_phase, clkout5_divide_f,
                           clkout5_duty_cycle, clkout5_phase, clkout6_divide_f,
                           clkout6_duty_cycle, clkout6_phase, ref_jitter1, ref_jitter2}),
                   str(compensation)
                   clock({clkfb, clkin1, clkin2, dclk, psclk}),
                   value({clkinsel, den, dwe, psen, psincdec, pwrdwn, reset}, "log"),
                   value(din, "bits:16"), value(daddr, "uint:7")
               ) {}
            \end{lstlisting}

            {\bf Output parameters:} \\
            \begin{tabular}{| l | l | l | l |}
            \hline
                               & mode  & type   & \\
            \hline \hline
            {\bf clkout0}      & clock &         & output clock 0 \\ \hline
            {\bf clkout1}      & clock &         & output clock 1 \\ \hline
            {\bf clkout2}      & clock &         & output clock 2 \\ \hline
            {\bf clkout3}      & clock &         & output clock 3 \\ \hline
            {\bf clkout4}      & clock &         & output clock 4 \\ \hline
            {\bf clkout5}      & clock &         & output clock 5 \\ \hline
            {\bf clkout6}      & clock &         & output clock 6 \\ \hline
            {\bf clkout0b}     & clock &         & output inverted clock 0 \\ \hline
            {\bf clkout1b}     & clock &         & output inverted clock 1 \\ \hline
            {\bf clkout2b}     & clock &         & output inverted clock 2 \\ \hline
            {\bf clkout3b}     & clock &         & output inverted clock 3 \\ \hline
            {\bf clkfbout}     & clock &         & output clock feedback \\ \hline
            {\bf clkfboutb}    & clock &         & output inverted clock feedback \\ \hline
            {\bf clkfbstopped} & value & log     & clock feedback stopped indicator \\ \hline
            {\bf clkinstopped} & value & log     & input clock stopped indicator \\ \hline
            {\bf drdy}         & value & log     & DRP ready \\ \hline
            {\bf locked}       & value & log     & MMC locked \\ \hline
            {\bf psdone}       & value & log     & phase shift done \\ \hline
            {\bf dout}         & value & log     & DRP data output \\ \hline
            \end{tabular}

            {\bf Input parameters:} \\            
            \begin{tabular}{| l | l | l | l |}
            \hline
                                          & mode  & type   & \\
            \hline \hline
            {\bf loc}                     & immediate & str   & optional location \\ \hline
            {\bf clkfbout\_use\_fine\_ps} & immediate & log   & enable feedback phase shift fine counter \\ \hline
            {\bf clkout0\_use\_fine\_ps}  & immediate & log   & enable output 0 phase shift fine counter  \\ \hline
            {\bf clkout1\_use\_fine\_ps}  & immediate & log   & enable output 1 phase shift fine counter  \\ \hline
            {\bf clkout2\_use\_fine\_ps}  & immediate & log   & enable output 2 phase shift fine counter  \\ \hline
            {\bf clkout3\_use\_fine\_ps}  & immediate & log   & enable output 3 phase shift fine counter  \\ \hline
            {\bf clkout4\_use\_fine\_ps}  & immediate & log   & enable output 4 phase shift fine counter  \\ \hline
            {\bf clkout5\_use\_fine\_ps}  & immediate & log   & enable output 5 phase shift fine counter  \\ \hline
            {\bf clkout6\_use\_fine\_ps}  & immediate & log   & enable output 6 phase shift fine counter  \\ \hline
            {\bf clock\_hold}             & immediate & log   & not documented \\ \hline
            {\bf startup\_wait}           & immediate & log   & delay startup on lock \\ \hline
            {\bf divclk\_divide}          & immediate & int   & divider ratio, 1 to 128 \\ \hline
            {\bf clkfbout\_phase}         & immediate & float & clock feedback phase, -360.0 to 360.0\\ \hline
            {\bf clkfbout\_mult\_f}       & immediate & float & clock feedback multiplier, 1.0 to 64.0\\ \hline
            {\bf clkout0\_divide\_f}      & immediate & float & clock 0 divider, 1.0 to 128.0\\ \hline
            {\bf clkout0\_duty\_cycle}    & immediate & float & clock 0 duty cycle, 0.001 to 0.999\\ \hline
            {\bf clkout0\_phase}          & immediate & float & clock 0 phase, -360.0 to 360.0\\ \hline
            {\bf clkout1\_divide\_f}      & immediate & float & clock 1 divider, 1.0 to 128.0\\ \hline
            {\bf clkout1\_duty\_cycle}    & immediate & float & clock 1 duty cycle, 0.001 to 0.999\\ \hline
            {\bf clkout1\_phase}          & immediate & float & clock 1 phase, -360.0 to 360.0\\ \hline
            {\bf clkout2\_divide\_f}      & immediate & float & clock 2 divider, 1.0 to 128.0\\ \hline
            {\bf clkout2\_duty\_cycle}    & immediate & float & clock 2 duty cycle, 0.001 to 0.999\\ \hline
            {\bf clkout2\_phase}          & immediate & float & clock 2 phase, -360.0 to 360.0\\ \hline
            {\bf clkout3\_divide\_f}      & immediate & float & clock 3 divider, 1.0 to 128.0\\ \hline
            {\bf clkout3\_duty\_cycle}    & immediate & float & clock 3 duty cycle, 0.001 to 0.999\\ \hline
            {\bf clkout3\_phase}          & immediate & float & clock 3 phase, -360.0 to 360.0\\ \hline
            {\bf clkout4\_divide\_f}      & immediate & float & clock 4 divider, 1.0 to 128.0\\ \hline
            {\bf clkout4\_duty\_cycle}    & immediate & float & clock 4 duty cycle, 0.001 to 0.999\\ \hline
            {\bf clkout4\_phase}          & immediate & float & clock 4 phase, -360.0 to 360.0\\ \hline
            {\bf clkout5\_divide\_f}      & immediate & float & clock 5 divider, 1.0 to 128.0\\ \hline
            {\bf clkout5\_duty\_cycle}    & immediate & float & clock 5 duty cycle, 0.001 to 0.999\\ \hline
            {\bf clkout5\_phase}          & immediate & float & clock 5 phase, -360.0 to 360.0\\ \hline
            {\bf clkout6\_divide\_f}      & immediate & float & clock 6 divider, 1.0 to 128.0\\ \hline
            {\bf clkout6\_duty\_cycle}    & immediate & float & clock 6 duty cycle, 0.001 to 0.999\\ \hline
            {\bf clkout6\_phase}          & immediate & float & clock 6 phase, -360.0 to 360.0\\ \hline
            {\bf ref\_jitter1}            & immediate & float & input clock 1 jitter, 0.000 to 0.999\\ \hline
            {\bf ref\_jitter2}            & immediate & float & input clock 2 jitter, 0.000 to 0.999\\ \hline
            {\bf compensation}            & immediate & str   & see Xilinx documentation\\ \hline
            {\bf clkfb}                   & clock     &       & clock feedback input   \\ \hline
            {\bf clkin1}                  & clock     &       & input clock 1   \\ \hline 
            {\bf clkin2}                  & clock     &       & input clock 2    \\ \hline
            {\bf dclk}                    & clock     &       & DRP clock input   \\ \hline
            {\bf psclk}                   & clock     &       & phase shift clock   \\ \hline 
            {\bf clkinsel}                & value     & log   & input clock selector \\ \hline
            {\bf den}                     & value     & log   & DRP enable  \\ \hline
            {\bf dwe}                     & value     & log   & DRP write enable   \\ \hline
            {\bf psen}                    & value     & log   & phase shift enable  \\ \hline
            {\bf psincdec}                & value     & log   & phase shift increment or decrement \\ \hline
            {\bf pwrdwn}                  & value     & log   & power down  \\ \hline
            {\bf reset}                   & value     & log   & reset  \\ \hline
            {\bf din}                     & value     & bits: & DRP data input   \\ \hline
            {\bf daddr}                   & value     & uint: & DRP address input    \\ \hline 
            \end{tabular}

            {\bf In-line target execution:} \\
            no

            {\bf description:} \\
            Creates an advanced Virtex6 or series 7 mixed mode clock manager. See Xilinx
            documentation for signal and attribute details. The first group of input
            parameters are immediate mode and are attributes. The remaining input
            parameters are input signals, clock mode for clocks and value mode for
            control signals.

        \subsubsection{clkmmcm\_base - create a basic MMCM multi-mode clock manager }\label{sec:lpmmcmb}
            
            {\bf Library: xilinx/clocks.3pl}

            {\bf Prototype:} \\
            \vsp
            \begin{lstlisting}[gobble=12, language=threepl]
               global procedure
               clkmmcm_base (
                   clock({ clkout0, clkout1, clkout2, clkout3, clkout4, clkout5, clkout6,
                           clkout0b, clkout1b, clkout2b, clkout3b, clkfbout, clkfboutb}),
                   value(locked, "log")
                   <-
                   str(loc),
                   str(bandwidth),
                   float(clkfbout_mult_f),
                   float(clkfbout_phase),
                   float(clkin1_period),
                   int({clkout1_divide, clkout2_divide, clkout3_divide,
                        clkout4_divide, clkout5_divide, clkout6_divide}),
                   float(clkout0_divide_f),
                   float(clkout0_duty_cycle, clkout1_duty_cycle,
                         clkout2_duty_cycle, clkout3_duty_cycle,
                         clkout4_duty_cycle, clkout5_duty_cycle,
                         clkout6_duty_cycle),
                   float(clkout0_phase, clkout1_phase, clkout2_phase,
                         clkout3_phase, clkout4_phase, clkout5_phase,
                         clkout6_phase),
                   int(divclk_divide),
                   float(ref_jitter1),
                   log(startup_wait),
                   clock({clkin, clkfb}),
                   value({pwrdwn, reset}, "log")
               ) {}
            \end{lstlisting}

            {\bf Output parameters:} \\
            \begin{tabular}{| l | l | l | l |}
            \hline
                               & mode  & type   & \\
            \hline \hline
            {\bf clkout0}      & clock &         & output clock 0 \\ \hline
            {\bf clkout1}      & clock &         & output clock 1 \\ \hline
            {\bf clkout2}      & clock &         & output clock 2 \\ \hline
            {\bf clkout3}      & clock &         & output clock 3 \\ \hline
            {\bf clkout4}      & clock &         & output clock 4 \\ \hline
            {\bf clkout5}      & clock &         & output clock 5 \\ \hline
            {\bf clkout6}      & clock &         & output clock 6 \\ \hline
            {\bf clkout0b}     & clock &         & output inverted clock 0 \\ \hline
            {\bf clkout1b}     & clock &         & output inverted clock 1 \\ \hline
            {\bf clkout2b}     & clock &         & output inverted clock 2 \\ \hline
            {\bf clkout3b}     & clock &         & output inverted clock 3 \\ \hline
            {\bf clkfbout}     & clock &         & output clock feedback \\ \hline
            {\bf clkfboutb}    & clock &         & output inverted clock feedback \\ \hline
            \end{tabular}

            {\bf Input parameters:} \\            
            \begin{tabular}{| l | l | l | l |}
            \hline
                                          & mode  & type   & \\
            \hline \hline
            {\bf loc}                     & immediate & str   & optionallocation \\ \hline
            {\bf bandwidth}               & immediate & str   & bandwidth \\ \hline
            {\bf clkfbout\_mult\_f}       & immediate & float & output multiplier \\ \hline
            {\bf clkfbout\_phase}         & immediate & float & output phase \\ \hline
            {\bf clkin1\_period}          & immediate & float & input clock period \\ \hline
            {\bf clkout0\_divide\_f}      & immediate & float & clock 0 divider, 1.0 to 128.0\\ \hline
            {\bf clkout1\_divide}         & immediate & int   & clock 1 divider, 1 to 128\\ \hline
            {\bf clkout2\_divide}         & immediate & int   & clock 2 divider, 1 to 128\\ \hline
            {\bf clkout3\_divide}         & immediate & int   & clock 3 divider, 1 to 128\\ \hline
            {\bf clkout4\_divide}         & immediate & int   & clock 4 divider, 1 to 128\\ \hline
            {\bf clkout5\_divide}         & immediate & int   & clock 5 divider, 1 to 128\\ \hline
            {\bf clkout6\_divide}         & immediate & int   & clock 6 divider, 1 to 128\\ \hline
            {\bf clkout0\_duty\_cycle}    & immediate & float & clock 0 duty cycle, 0.001 to 0.999\\ \hline
            {\bf clkout1\_duty\_cycle}    & immediate & float & clock 1 duty cycle, 0.001 to 0.999\\ \hline
            {\bf clkout2\_duty\_cycle}    & immediate & float & clock 2 duty cycle, 0.001 to 0.999\\ \hline
            {\bf clkout3\_duty\_cycle}    & immediate & float & clock 3 duty cycle, 0.001 to 0.999\\ \hline
            {\bf clkout4\_duty\_cycle}    & immediate & float & clock 4 duty cycle, 0.001 to 0.999\\ \hline
            {\bf clkout5\_duty\_cycle}    & immediate & float & clock 5 duty cycle, 0.001 to 0.999\\ \hline
            {\bf clkout6\_duty\_cycle}    & immediate & float & clock 6 duty cycle, 0.001 to 0.999\\ \hline
            {\bf clkout0\_phase}          & immediate & float & clock 0 phase, -360.0 to 360.0\\ \hline
            {\bf clkout1\_phase}          & immediate & float & clock 1 phase, -360.0 to 360.0\\ \hline
            {\bf clkout2\_phase}          & immediate & float & clock 2 phase, -360.0 to 360.0\\ \hline
            {\bf clkout3\_phase}          & immediate & float & clock 3 phase, -360.0 to 360.0\\ \hline
            {\bf clkout4\_phase}          & immediate & float & clock 4 phase, -360.0 to 360.0\\ \hline
            {\bf clkout5\_phase}          & immediate & float & clock 5 phase, -360.0 to 360.0\\ \hline
            {\bf clkout6\_phase}          & immediate & float & clock 6 phase, -360.0 to 360.0\\ \hline
            {\bf divclk\_divide}          & immediate & int   & divider ratio, 1 to 128 \\ \hline
            {\bf startup\_wait}           & immediate & log   & delay startup on lock \\ \hline
            {\bf clkin}                   & clock     &       & input clock   \\ \hline 
            {\bf clkfb}                   & clock     &       & clock feedback input   \\ \hline
            {\bf pwrdwn}                  & value     & log   & power down  \\ \hline
            {\bf reset}                   & value     & log   & reset  \\ \hline
            \end{tabular}

            {\bf In-line target execution:} \\
            no

            {\bf description:} \\
            Creates a basic Virtex6 or series 7 mixed mode clock manager. See Xilinx
            documentation for signal and attribute details. The first group of input
            parameters are immediate mode and are attributes. The remaining input
            parameters are input signals, clock mode for clocks and value mode for
            control signals.



        \subsubsection{pll\_base - create a series-7 basic PLL}\label{sec:lmpllb}
            
            {\bf Library: xilinx/clocks.3pl}

            {\bf Prototype:} \\
            \vsp
            \begin{lstlisting}[gobble=12, language=threepl]
               global procedure
               pll_base (
                   clock({clkout0, clkout1, clkout2, clkout3, clkout4, clkout5, clkfbout}),
                   value(locked, "log")
                   <-
                   str(loc),
                   int({   clkout0_divide, clkout1_divide, clkout2_divide, clkout3_divide,
                           clkout4_divide, clkout5_divide, clkfbout_mult, clkfbout_phase,
                           divclk_divide}),
                   str({compensation, bandwidth}),
                   float({clkout0_phase, clkout1_phase, clkout2_phase, clkout3_phase, clkout4_phase,
                           clkout5_phase, clkout0_duty_cycle, clkout1_duty_cycle, clkout2_duty_cycle,
                           clkout3_duty_cycle, clkout4_duty_cycle, clkout5_duty_cycle, ref_jitter,
                           clkin_period}),
                   clock({clkin, clkfbin}), value(reset, "log")
               ) {}
            \end{lstlisting}

            {\bf Output parameters:} \\
            \begin{tabular}{| l | l | l | l |}
            \hline
                           & mode  & type   & \\
            \hline \hline
            {\bf clkout0}  & clock &         & output clock 0 \\ \hline    
            {\bf clkout1}  & clock &         & output clock 0 \\ \hline    
            {\bf clkout2}  & clock &         & output clock 0 \\ \hline    
            {\bf clkout3}  & clock &         & output clock 0 \\ \hline    
            {\bf clkout4}  & clock &         & output clock 0 \\ \hline    
            {\bf clkout5}  & clock &         & output clock 0 \\ \hline    
            {\bf clkfbout} & clock &         & feedback output \\ \hline   
            {\bf locked}   & value & log     & PLL is locked \\ \hline
            \end{tabular}

            {\bf Input parameters:} \\            
            \begin{tabular}{| l | l | l | l |}
            \hline
                           & mode  & type   & \\
            \hline \hline
            {\bf loc}                  & immediate & str   & optional location \\ \hline
            {\bf clkout0\_divide}      & immediate & int   & division ratio for output 0\\ \hline
            {\bf clkout1\_divide}      & immediate & int   & division ratio for output 1\\ \hline
            {\bf clkout2\_divide}      & immediate & int   & division ratio for output 2\\ \hline
            {\bf clkout3\_divide}      & immediate & int   & division ratio for output 3\\ \hline
            {\bf clkout4\_divide}      & immediate & int   & division ratio for output 4\\ \hline
            {\bf clkout5\_divide}      & immediate & int   & division ratio for output 5\\ \hline
            {\bf clkfbout\_mult}       & immediate & int   & multiplication factor for all outputs \\ \hline
            {\bf clkfbout\_phase}      & immediate & int   & phase shift for feedback \\ \hline
            {\bf divclk\_divide}       & immediate & int   & division ration for all outputs \\ \hline
            {\bf compensation}         & immediate & str   &             \\ \hline
            {\bf bandwidth}            & immediate & str   &             \\ \hline
            {\bf clkout0\_phase}       & immediate & float & phase shift for output 0 \\ \hline
            {\bf clkout1\_phase}       & immediate & float & phase shift for output 1 \\ \hline
            {\bf clkout2\_phase}       & immediate & float & phase shift for output 2 \\ \hline
            {\bf clkout3\_phase}       & immediate & float & phase shift for output 3 \\ \hline
            {\bf clkout4\_phase}       & immediate & float & phase shift for output 4 \\ \hline
            {\bf clkout5\_phase}       & immediate & float & phase shift for output 5 \\ \hline
            {\bf clkout0\_duty\_cycle} & immediate & float & duty cycle for output 0 \\ \hline
            {\bf clkout1\_duty\_cycle} & immediate & float & duty cycle for output 1 \\ \hline
            {\bf clkout2\_duty\_cycle} & immediate & float & duty cycle for output 2 \\ \hline
            {\bf clkout3\_duty\_cycle} & immediate & float & duty cycle for output 3 \\ \hline
            {\bf clkout4\_duty\_cycle} & immediate & float & duty cycle for output 4 \\ \hline
            {\bf clkout5\_duty\_cycle} & immediate & float & duty cycle for output 5 \\ \hline
            {\bf ref\_jitter}          & immediate & float & reference clock jitter \\ \hline
            {\bf clkin\_period}        & immediate & float & input clock period, ns\\ \hline
            {\bf clkin}                & clock     &       & input clock \\ \hline
            {\bf clkfbin}              & clock     &       & clock feedback \\ \hline
            {\bf reset}                & value     & log   & reset \\ \hline
            \end{tabular}

            {\bf In-line target execution:} \\
            no

            {\bf description:} \\
            Creates a Virtex5 PLL phase locked loop. See Xilinx documentation for
            signal and attribute details. The first group of input parameters are
            immediate mode and are attributes. The remaining input parameters are
            input signals, clock mode for clocks and value mode for control signals.

            The PLL is PLL\_BASE, i.e. only a single input is provided.
            The dual switchable input PLL, PLL\_ADV, is not supported.

            \heading{attributes}

                {\bf loc} is an optional location string from the
                following list - "DLL0P", "DLL1P", "DLL2P", "DLL3P", "DLL0S", "DLL1S",
                "DLL2S", "DLL3S". These select from the four primary and four secondary
                delay-locked loops. Ignored if "".

                {\bf clkin\_period} is the input clock period. This must be
                specified if the DFS is used.

                {\bf compensation} specifies the PLL phase compensation. The values are
                "system\_synchronous" (default) or "source\_synchronous".

                {\bf bandwidth} specifies the PLL programming algorithm affecting jitter,
                phase margin and other characteristics. Values are "high", "low" and
                "optimised" (default).

                {\bf clkout0\_divide} to {\bf clkout5\_divide} provides a PLL division ratio
                for each clock output. The values range from 1 (default) to 128.

                {\bf clkout0\_phase} to {\bf clkout5\_phase} specifies the PLL output phase
                relationship (relative to what?). Values are in the range 0.01 to 360.0 but
                the default is listed as 0.0.

                {\bf clkout0\_duty\_cycle} to {\bf clkout5\_duty\_cycle} specifies the PLL
                duty cycle of each clock output. Vales range from 0.01 to 0.99 with a
                default of 0.5.

                {\bf clkfbout\_mult} specifies a common frequency multiplier for all PLL
                clock outputs. This works in combination with the individual frequency
                divisors. The values range from 1 to 64 with a default of 1.

                {\bf divclk\_divide} specifies a frequency division ration for all PLL clock
                outputs. It is documented how this interacts with the individual divisors.
                Range is 1 to 52 with a default of 1.

                {\bf clkfbout\_phase} specifies a phase offset for the PLL clock feedback
                output. Values range from 0.0 to 360.0 with a default of 0.0.

                {\bf ref\_jitter} specifies the PLL reference clock jitter. This is documented
                as ``specified interms of the UI which is a percentage of the reference
                clock''. Allowed values are 0.000 to 0.999 with a default of 0.100.


            \heading{input signals}

            {\bf clkin} is the input clock.

            {\bf clkfbin} is the feedback clock input, which is normally conected to
            {\bf clkfbout} or a buffered version of an output clock.

            {\bf reset} is the reset input.



            \heading{output signals}

                {\bf clkout0}, {\bf clkout1}, {\bf clkout2}, {\bf clkout3}, {\bf
                clkout4} and {\bf clkout5} are separate clock outputs for a PLL. Each
                of these has its own frequency divisor, phase and duty cycle.

                {\bf clkfbout} is the clock feedback output.

                {\bf locked} is high when the DLL and DFS are locked to the input
                clock. This signal has a null clock domain.


        \subsubsection{clockbufg - create a BUFG clock buffer}\label{sec:lpclockbufg}
            
            {\bf Library: xilinx/clocks.3pl}

            {\bf Prototype:} \\
            \vsp
            \begin{lstlisting}[gobble=12, language=threepl]
               global procedure
               clockbufg (clock(out) <- clock(in), str(loc), log(invert)) {
            \end{lstlisting}

            {\bf Output parameters:} \\
            {\bf out} is the clock variable for the distributed clock.

            {\bf Input parameters:} \\            
            {\bf in} is the internal clock variable for the buffer input
            clock.\\
            {\bf loc} is optional and is specified by a string of the form
            "BUFGMUX2S". Generally the location argument may be omitted or "",
            allowing the Xilinx software to select a clock buffer. \\
            {\bf invert} specifies an optional inverter. Ignored if "".
        
            {\bf In-line target execution:} \\
            no

            {\bf description:} \\
            This procedure creates a BUFG clock buffer which is common to all
            Xilinx FPGAs. The number of clock buffers
            available depends on the particular FPGA.
            
            If {\bf invert} is true an inverter is inserted at the input.
            
            The frequency, minimum frequency, maximum frequency, period and duty
            cycle of the output clock are derived from the input clock overriding
            any previous values.



        \subsubsection{clockbufio - a local clock buffer}\label{sec:lpclockbufio}
            
            {\bf Library: xilinx/clocks.3pl}

            {\bf Prototype:} \\
            \vsp
            \begin{lstlisting}[gobble=12, language=threepl]
               global procedure
               clockbufio (clock(out) <- clock(in), str(loc)) {}
            \end{lstlisting}

            {\bf Output parameters:} \\
            {\bf out} is the clock variable for the distributed clock.

            {\bf Input parameters:} \\            
            {\bf in} is the internal clock variable for the buffer input
            clock.\\
            {\bf loc} is optional and is specified by a string of the form
            "BUFGMUX2S". Generally the location argument may be omitted or "",
            allowing the Xilinx software to select a clock buffer. \\
        
            {\bf In-line target execution:} \\
            no

            {\bf description:} \\
            This procedure creates a local clock buffer. The number of clock
            buffers available depends on the particular FPGA.
            
            The frequency, period and duty cycle of the output clock are
            derived from the input clock overriding any previous values.



        \subsubsection{clockmuxbuffer - a clock mux/buffer}\label{sec:lpclockmuxbuffer}
            
            {\bf Library: xilinx/clocks.3pl}

            {\bf Prototype:} \\
            \vsp
            \begin{lstlisting}[gobble=12, language=threepl]
               global procedure
               clockmuxbuffer (clock(out) <- clock({in0, in1}), value(select, "log"), str(loc) {}
            \end{lstlisting}

            {\bf Output parameters:} \\
            {\bf out} is the clock variable for the distributed clock.

            {\bf Input parameters:} \\            
            {\bf in0} is the internal clock variable for the first buffer input
            clock.\\
            {\bf in1} is the internal clock variable for the second buffer input
            clock.\\
            {\bf select} is a logical value which selects the first clock
            input if false and the second clock input if true.\\
            {\bf loc} is optional and is specified by a string of the form
            "BUFGMUX2S". Generally the location argument may be omitted,
            allowing the Xilinx software to select a clock buffer.
        
            {\bf In-line target execution:} \\
            no

            {\bf description:} \\
            This procedure creates a clock buffer with multiplexed input. The
            number of multiplexed clock buffers available depends on the
            particular FPGA.
            
            The frequency, period and duty cycle of the output clock are
            derived from the input clock overriding any previous values.



        \subsubsection{clockmuxbufferdisable - a clock mux/buffer with disable}\label{sec:lpclockmuxbufferdis}
            
            {\bf Library: xilinx/clocks.3pl}

            {\bf Prototype:} \\
            \vsp
            \begin{lstlisting}[gobble=12, language=threepl]
               global procedure
               clockmuxbufferdisable (clock(out) <- clock(in), value(disable, "log"), str(loc)) {
            \end{lstlisting}

            {\bf Output parameters:} \\
            {\bf out} is the clock variable for the distributed clock.

            {\bf Input parameters:} \\            
            {\bf in} is the internal clock variable for the first buffer input
            clock.\\
            {\bf disable} is a logical value which disables the clock
            output if true.\\
            {\bf loc} is optional and is specified by a string of the form
            "BUFGMUX2S". Generally the location argument may be omitted or "",
            allowing the Xilinx software to select a clock buffer.
        
            {\bf In-line target execution:} \\
            no

            {\bf description:} \\
            This procedure creates a clock buffer with enabled output. If the
            argument to parameter {\bf disable} goes true the clock output
            is disabled low. The output clock is enabled and disabled without
            glitches or shortened high half-periods.
            
            The frequency, period and duty cycle of the output clock are
            derived from the input clock overriding any previous values.




        \subsubsection{clock\_params - derive clock synthesis parameters}\label{sec:lpclockparams}
            
            {\bf Library: xilinx/clocks.3pl}

            {\bf Prototype:} \\
            \vsp
            \begin{lstlisting}[gobble=12, language=threepl]
               global procedure
               clock_params (
                   int({mul, div}),
                   float(af)
                   <-
                   float({fin, freq}),
                   int({mmin, mmax, dmin, dmax}),
                   float(tol)
               ) {}
            \end{lstlisting}
            
            {\bf Output parameters:} \\
            {\bf mul} is the synthesis multiplier \\
            {\bf div} is the synthesis divisor \\
            {\bf af} is the resulting frequency

            {\bf Input parameters:} \\
            {\bf fin} is the input frequency \\
            {\bf freq} is the required frequency \\
            {\bf mmin} is the minimum value of the multiplier \\
            {\bf mmax} is the maximum value of the multiplier \\
            {\bf dmin} is the minimum value of the divisor \\
            {\bf dmax} is the maximum value of the divisor \\
            {\bf tol} is an optional frequency tolerance to minimise the search
        
            {\bf In-line target execution:} \\
            no

            {\bf description:} \\
            This procedure derives an integer multiplier and an integer divisor
            for clock synthesis.

            The procedure searches through combinations of multiplier and divisor
            to find the nearest generated frequency that is not greater than the
            requested frequency. The optional tolerance argument when supplied
            greatly reduces the execution time. Frequencies may be in any units,
            e.g. MHz, KHz, Hz, since only ratios are being examined, but hence the
            units must be consistent. It is nevertheless preferable to use
            frequencies in Hz to be consistent in all contexts.
            
            If the requested frequency cannot be generated within the
            contraints of the arguments supplied the procedure returns
            an error code.
            
            {\bf mul} and {\bf div} are the frequency synthesis parameters,
            {\bf af} is the actual frequency achieved.
            
            A fatal error occurs if the frequency cannot be synthesised because the tolerance is too
            small or the requested frequency is too low (the error message will specify which).
            
            To avoid the fatal error the related library function of the same name can be used
            instead. The funstion returns a structure containing the result parameters along
            with an error code.
            
            

        \subsubsection{lowskew - request low skew secondary clock routing}\label{sec:lplowskew}
            \index{constraint!timing}
            \index{timing constraint}
            
            {\bf Library: xilinx/clocks.3pl}

            {\bf Prototype:} \\
            \vsp
            \begin{lstlisting}[gobble=12, language=threepl]
               global procedure
               lowskew (clock(clk)) {}
            \end{lstlisting}

            {\bf Output parameters:} \\
            none

            {\bf Input parameters:} \\            
            {\bf clk} is a clock variable.
        
            {\bf In-line target execution:} \\
            no

            {\bf description:} \\
            This procedure places a low skew constraint on the routing
            of a clock signal. This is required if the clock cannot be
            distributed by a global clock net (from a global clock buffer).
            There are only 4 global clock buffers.





    \subsection{library functions}

        \subsubsection{clock\_params - derive clock synthesis parameters}\label{sec:lfclockparams}
            
            {\bf Library: xilinx/clocks.3pl}

            {\bf Prototype:} \\
            \vsp
            \begin{lstlisting}[gobble=12, language=threepl]
               global function
               clock_params (
                   float({fin, freq}),
                   int({mmin, mmax, dmin, dmax}),
                   float(tol)
               ) {}
            \end{lstlisting}

            {\bf Input parameters:} \\            
            {\bf fin} is the input frequency \\
            {\bf freq} is the required frequency \\
            {\bf mmin} is the minimum value of the multiplier \\
            {\bf mmax} is the maximum value of the multiplier \\
            {\bf dmin} is the minimum value of the divisor \\
            {\bf dmax} is the maximum value of the divisor \\
            {\bf tol} is an optional frequency tolerance to minimise the search

            {\bf Return mode and type:} \\
            immediate, struct
        
            {\bf In-line target execution:} \\
            no

            {\bf description:} \\
            This function derives an integer multiplier and an integer divisor
            for clock synthesis.

            The function searches through combinations of multiplier and divisor
            to find the nearest generated frequency that is not greater than the
            requested frequency. The optional tolerance argument when supplied
            greatly reduces the execution time. Frequencies may be in any units,
            e.g. MHz, KHz, Hz, since only ratios are being examined, but hence the
            units must be consistent. It is nevertheless preferable to use
            frequencies in Hz to be consistent in all contexts.
            
            
            If the requested frequency cannot be generated within the
            contraints of the arguments supplied the procedure returns
            an error code.
            
            The return value has the following structure -
            
            \begin{lstlisting}[gobble=12, language=threepl]
               struct(rv,
                   {mul, div}, "int",
                   af,         "float",
                   exit,       "int"
               );
            \end{lstlisting}
            where {\bf mul} and {\bf div} are the frequency synthesis parameters,
            {\bf af} is the actual frequency achieved and {\bf exit} is an exit code
            which is 0 for normal return, 1 for a frequency tolerance which is too
            small and 2 for a requested frequency which is too low.
            



    \subsection{library classes}
    
        None.

